\subsection{\textit{Dynamic Difficulty Ajustment}}
    O \textit{Dynamic Difficulty Adjustment} (DDA) 
    consiste em duas etapas \cite{zohaib2018dynamic}.
    Primeiramente, mede-se o nível de habilidade dos 
    jogadores por meio de uma função de pontuação. 
    Em seguida, ajusta-se a dificuldade do jogo de 
    acordo com os níveis de habilidade estimados, de 
    diversas formas, como o ajuste de parâmetros no jogo
    \cite{denisova2015adaptation},
    a manipulação do ambiente de jogo
    \cite{duque2020finding,lora2016dynamic} 
    e a utilização de agentes de jogo
    \cite{demediuk2017monte,ishihara2018monte}.